\chapter*{Appendix C: Architecture Benchmark Data}
\label{appendix:supplementary}
\addcontentsline{toc}{chapter}{Appendix C: Architecture Benchmark Data}

\section*{C.1 Architecture Benchmark Detail}
\label{appendix:benchmark}

\subsection*{C.1.1 In-process LRU cache counterfactual}
\label{sec:solution:architecture:lru}

A natural objection to materialising a Gold tier is that the same
latency could be achieved by an in-process least-recently-used (LRU)
cache wrapped around the Silver-only query. The counterfactual is
worth measuring rather than dismissing on principle: a process-local
LRU is simpler than a persistent Gold table and incurs no warehouse
build step. Table~\ref{tab:arch:lru} measures it against the same six
aggregates of Table~\ref{tab:arch:quantitative}, with cold and warm
runs separated.

\begin{table}[!htbp]
\centering
\caption{LRU result cache versus persistent Gold tables, cold and warm runs.}
\label{tab:arch:lru}
\footnotesize
\renewcommand{\arraystretch}{1.20}
\begin{tabularx}{\linewidth}{p{4.8cm} r r r X}
\hline
\textbf{Aggregate} & \textbf{LRU cold} & \textbf{LRU warm} & \textbf{Gold} & \textbf{Outcome} \\
\hline
\texttt{agg\_network\_day}     &   893.7 ms &  $<$ 0.1 ms &    0.4 ms & Gold $\ll$ LRU cold \\
\texttt{agg\_route\_hour}      & 1{,}324.7 ms &  $<$ 0.1 ms &   35.0 ms & Gold $\ll$ LRU cold \\
\texttt{agg\_stop\_hour}       & 2{,}453.7 ms &  $<$ 0.1 ms &  852.7 ms & Gold $\ll$ LRU cold \\
\texttt{agg\_school\_route}    &   319.3 ms &  $<$ 0.1 ms &   30.3 ms & Gold $\ll$ LRU cold \\
\texttt{agg\_stop\_total}      &   477.2 ms &  $<$ 0.1 ms &    4.7 ms & Gold $\ll$ LRU cold \\
\texttt{agg\_hub\_kpis}        &   272.1 ms &  $<$ 0.1 ms &    0.3 ms & Gold $\ll$ LRU cold \\
\hline
\end{tabularx}
\end{table}

The cache-warm path is sub-millisecond and beats the Gold read on
every aggregate, which would be a point in favour of the cache if
warmth were free. The cache-cold path costs exactly what the
Silver-only mode of Table~\ref{tab:arch:quantitative} costs: the
first call against an empty cache executes the full Silver-only
query, and a process restart resets every entry. The cache also has
two further weaknesses that a Gold table does not: a per-process
memory ceiling that grows with cache entries times result-set size,
and no cross-process sharing (a multi-user deployment with $k$
worker processes pays the cold cost $k$ times). The Gold tier is
persistent on disk, shared across every worker process, and reload
on server start is bounded at the size of the materialised table
($\sim$30~MB for the seven Gold aggregates combined).

\subsection*{C.1.2 Build-time scaling validation}
\label{sec:solution:architecture:scaling}

The build-time penalty of materialising the Gold tier is paid offline
at data-refresh boundaries, but the penalty must scale predictably
against the underlying input volume for the architecture to remain
defensible as the data grows. Table~\ref{tab:arch:scaling} validates
the scaling empirically by running the heaviest single Gold build
(\texttt{agg\_route\_hour}) at one and two times the baseline input
size and measuring the build-time ratio.

\begin{table}[!htbp]
\centering
\caption{Build-time scaling for \texttt{agg\_route\_hour}, median of three runs.}
\label{tab:arch:scaling}
\footnotesize
\renewcommand{\arraystretch}{1.20}
\begin{tabularx}{\linewidth}{p{5.8cm} r X}
\hline
\textbf{Input size} & \textbf{Build time} & \textbf{Notes} \\
\hline
1$\times$ input ($17{,}883{,}623$ rows) & 1{,}361.8 ms & Measured. \\
2$\times$ input ($35{,}767{,}246$ rows) & 2{,}650.2 ms & Measured. \\
Empirical scaling ratio                 & \textbf{1.95$\times$} & Ideal O($N$) = 2.00$\times$; 2.5\% deviation is fixed-overhead outside the row scan. \\
5$\times$ input ($\approx$89.4~M rows)  & $\approx$6{,}809 ms & Linear projection. \\
10$\times$ input ($\approx$178.8~M rows) & $\approx$13{,}618 ms & Linear projection. \\
\hline
\end{tabularx}
\end{table}

The empirical scaling ratio of 1.95$\times$ for a 2$\times$ input
confirms that the Gold build-time cost is linear in input volume,
with a small constant overhead that becomes proportionally smaller
at larger scales.
